% =====================================================================
%  VERIFICA DI SISTEMI AUTOMATICI 3BETC — dati ed esercizi.
%  Per ottenere i PDF imposta come Main document uno dei file numerati
%  di questa cartella:
%     1_Testo_FilaA   2_Testo_FilaB   3_Testo_FilaC          (da stampare)
%     4_Soluzioni_FilaA  5_Soluzioni_FilaB  6_Soluzioni_FilaC  (da dare agli alunni)
%     7_Correzione                                            (griglia per il docente)
%  Struttura (dagli appunti di Alessio del 20/10/2026):
%     1) decimale -> binario   2) binario -> decimale
%     3-4) 3 ingressi, 2 porte logiche in cascata
%     5-6) diagrammi di flusso A SCELTA: 5 facile (6 pt) oppure 6 difficile (14 pt)
% =====================================================================
% valori predefiniti se compili direttamente questo file
\providecommand{\fila}{A}          % A | B | C
\providecommand{\versione}{testo}  % testo | soluzioni | correzione | completa

\documentclass[a4paper,11pt,addpoints]{exam}
\usepackage{stili/verifica}

% ---------- Dati della prova --------------------------------------------
\scuola{ITIS G. Galilei -- Livorno}
\annoscolastico{2026/27}
\tipoprova{Verifica}
\materia{Sistemi Automatici}
\argomento{Sistema binario, diagrammi di flusso, porte logiche}
\classe{3BETC}
\dataprova{23/10/2026}
\durata{}               % non indicata sul foglio
\votominimo{2}          % voto con 0 punti
\sogliasufficienza{60}  % % del punteggio per il 6
\criterigenerali{Conversioni: errore di calcolo isolato con procedimento corretto: metà dei punti.
  Porte logiche: colonna $Y$ 5 pt, $-1$ pt per ogni riga sbagliata, minimo 0 (non si va in negativo).
  Diagrammi di flusso: lo studente svolge \textbf{uno solo} tra l'esercizio 5 (più semplice, max 6 pt)
  e l'esercizio 6 (più difficile, max 14 pt). Punteggio massimo della prova: 42 pt con l'esercizio 6,
  34 pt con l'esercizio 5 (voto massimo circa 8). Se li svolge entrambi si conta quello con il
  punteggio più alto. Si valuta la logica dell'algoritmo; simboli sbagliati o frecce mancanti
  $-0{,}5$ pt (una sola volta); soluzioni diverse ma corrette prendono tutti i punti.}
\puntimassimi{42}       % 8+8+6+6 + max(6, 14): es. 5 e 6 a scelta
% tabella punti: es. 5 e 6 a scelta, nel totale conta solo uno dei due
\renewcommand{\tabellapunteggi}{%
  \begin{center}\small\renewcommand{\arraystretch}{1.5}%
  \begin{tabular}{|l|c|c|c|c|c|c|}\hline
    Esercizio & 1 & 2 & 3 & 4 & 5 \emph{oppure} 6 & Totale\\\hline
    Punti     & 8 & 8 & 6 & 6 & 6 \emph{oppure} 14 & 42\\\hline
    Ottenuti  &   &   &   &   &                   & \\\hline
  \end{tabular}\end{center}}

\begin{document}

\intestazione

\begin{questions}

% ---------------------------------------------------------------------
% Argomento: conversione decimale -> binario (divisioni successive per 2)
%   4 numeri a 3 cifre per fila (8, 9, 9 e 10 bit), un solo numero pari
%   fila A: 173=10101101, 461=111001101, 358=101100110, 829=1100111101
%   fila B: 394=110001010, 211=11010011, 917=1110010101, 437=110110101
%   fila C: 489=111101001, 157=10011101, 763=1011111011, 326=101000110
% Livello: base   |   Tempo stimato: 12 min
% ---------------------------------------------------------------------
\begin{esercizio}[8]
Converti in binario:
\ABC{%
\par\smallskip\noindent\hspace*{1em}\begin{tabular}{@{}p{0.42\linewidth}p{0.42\linewidth}@{}}
  \textbf{a)}~$173_{10} = \ldots\ldots\ldots_2$ & \textbf{b)}~$461_{10} = \ldots\ldots\ldots_2$\\[4pt]
  \textbf{c)}~$358_{10} = \ldots\ldots\ldots_2$ & \textbf{d)}~$829_{10} = \ldots\ldots\ldots_2$
\end{tabular}}{%
\par\smallskip\noindent\hspace*{1em}\begin{tabular}{@{}p{0.42\linewidth}p{0.42\linewidth}@{}}
  \textbf{a)}~$394_{10} = \ldots\ldots\ldots_2$ & \textbf{b)}~$211_{10} = \ldots\ldots\ldots_2$\\[4pt]
  \textbf{c)}~$917_{10} = \ldots\ldots\ldots_2$ & \textbf{d)}~$437_{10} = \ldots\ldots\ldots_2$
\end{tabular}}{%
\par\smallskip\noindent\hspace*{1em}\begin{tabular}{@{}p{0.42\linewidth}p{0.42\linewidth}@{}}
  \textbf{a)}~$489_{10} = \ldots\ldots\ldots_2$ & \textbf{b)}~$157_{10} = \ldots\ldots\ldots_2$\\[4pt]
  \textbf{c)}~$763_{10} = \ldots\ldots\ldots_2$ & \textbf{d)}~$326_{10} = \ldots\ldots\ldots_2$
\end{tabular}}

\begin{solution}
\ABC{%
\noindent
  \begin{minipage}[t]{0.24\linewidth}
  \textbf{a)}\\[2pt]
  {\small\renewcommand{\arraystretch}{1.0}\begin{tabular}{r|l}
    173 & 1\\
    86 & 0\\
    43 & 1\\
    21 & 1\\
    10 & 0\\
    5 & 1\\
    2 & 0\\
    1 & 1\\
    0 & \\
  \end{tabular}}
  \end{minipage}\hfill
  \begin{minipage}[t]{0.24\linewidth}
  \textbf{b)}\\[2pt]
  {\small\renewcommand{\arraystretch}{1.0}\begin{tabular}{r|l}
    461 & 1\\
    230 & 0\\
    115 & 1\\
    57 & 1\\
    28 & 0\\
    14 & 0\\
    7 & 1\\
    3 & 1\\
    1 & 1\\
    0 & \\
  \end{tabular}}
  \end{minipage}\hfill
  \begin{minipage}[t]{0.24\linewidth}
  \textbf{c)}\\[2pt]
  {\small\renewcommand{\arraystretch}{1.0}\begin{tabular}{r|l}
    358 & 0\\
    179 & 1\\
    89 & 1\\
    44 & 0\\
    22 & 0\\
    11 & 1\\
    5 & 1\\
    2 & 0\\
    1 & 1\\
    0 & \\
  \end{tabular}}
  \end{minipage}\hfill
  \begin{minipage}[t]{0.24\linewidth}
  \textbf{d)}\\[2pt]
  {\small\renewcommand{\arraystretch}{1.0}\begin{tabular}{r|l}
    829 & 1\\
    414 & 0\\
    207 & 1\\
    103 & 1\\
    51 & 1\\
    25 & 1\\
    12 & 0\\
    6 & 0\\
    3 & 1\\
    1 & 1\\
    0 & \\
  \end{tabular}}
  \end{minipage}\\[6pt]
\par\smallskip\noindent\hspace*{1em}\begin{tabular}{@{}p{0.42\linewidth}p{0.42\linewidth}@{}}
  \textbf{a)}~$173_{10} = 10101101_2$ & \textbf{b)}~$461_{10} = 111001101_2$\\[4pt]
  \textbf{c)}~$358_{10} = 101100110_2$ & \textbf{d)}~$829_{10} = 1100111101_2$
\end{tabular}\par\smallskip
  Si divide per 2 fino ad arrivare a 0 e si leggono i resti dal basso verso l'alto.}{%
\noindent
  \begin{minipage}[t]{0.24\linewidth}
  \textbf{a)}\\[2pt]
  {\small\renewcommand{\arraystretch}{1.0}\begin{tabular}{r|l}
    394 & 0\\
    197 & 1\\
    98 & 0\\
    49 & 1\\
    24 & 0\\
    12 & 0\\
    6 & 0\\
    3 & 1\\
    1 & 1\\
    0 & \\
  \end{tabular}}
  \end{minipage}\hfill
  \begin{minipage}[t]{0.24\linewidth}
  \textbf{b)}\\[2pt]
  {\small\renewcommand{\arraystretch}{1.0}\begin{tabular}{r|l}
    211 & 1\\
    105 & 1\\
    52 & 0\\
    26 & 0\\
    13 & 1\\
    6 & 0\\
    3 & 1\\
    1 & 1\\
    0 & \\
  \end{tabular}}
  \end{minipage}\hfill
  \begin{minipage}[t]{0.24\linewidth}
  \textbf{c)}\\[2pt]
  {\small\renewcommand{\arraystretch}{1.0}\begin{tabular}{r|l}
    917 & 1\\
    458 & 0\\
    229 & 1\\
    114 & 0\\
    57 & 1\\
    28 & 0\\
    14 & 0\\
    7 & 1\\
    3 & 1\\
    1 & 1\\
    0 & \\
  \end{tabular}}
  \end{minipage}\hfill
  \begin{minipage}[t]{0.24\linewidth}
  \textbf{d)}\\[2pt]
  {\small\renewcommand{\arraystretch}{1.0}\begin{tabular}{r|l}
    437 & 1\\
    218 & 0\\
    109 & 1\\
    54 & 0\\
    27 & 1\\
    13 & 1\\
    6 & 0\\
    3 & 1\\
    1 & 1\\
    0 & \\
  \end{tabular}}
  \end{minipage}\\[6pt]
\par\smallskip\noindent\hspace*{1em}\begin{tabular}{@{}p{0.42\linewidth}p{0.42\linewidth}@{}}
  \textbf{a)}~$394_{10} = 110001010_2$ & \textbf{b)}~$211_{10} = 11010011_2$\\[4pt]
  \textbf{c)}~$917_{10} = 1110010101_2$ & \textbf{d)}~$437_{10} = 110110101_2$
\end{tabular}\par\smallskip
  Si divide per 2 fino ad arrivare a 0 e si leggono i resti dal basso verso l'alto.}{%
\noindent
  \begin{minipage}[t]{0.24\linewidth}
  \textbf{a)}\\[2pt]
  {\small\renewcommand{\arraystretch}{1.0}\begin{tabular}{r|l}
    489 & 1\\
    244 & 0\\
    122 & 0\\
    61 & 1\\
    30 & 0\\
    15 & 1\\
    7 & 1\\
    3 & 1\\
    1 & 1\\
    0 & \\
  \end{tabular}}
  \end{minipage}\hfill
  \begin{minipage}[t]{0.24\linewidth}
  \textbf{b)}\\[2pt]
  {\small\renewcommand{\arraystretch}{1.0}\begin{tabular}{r|l}
    157 & 1\\
    78 & 0\\
    39 & 1\\
    19 & 1\\
    9 & 1\\
    4 & 0\\
    2 & 0\\
    1 & 1\\
    0 & \\
  \end{tabular}}
  \end{minipage}\hfill
  \begin{minipage}[t]{0.24\linewidth}
  \textbf{c)}\\[2pt]
  {\small\renewcommand{\arraystretch}{1.0}\begin{tabular}{r|l}
    763 & 1\\
    381 & 1\\
    190 & 0\\
    95 & 1\\
    47 & 1\\
    23 & 1\\
    11 & 1\\
    5 & 1\\
    2 & 0\\
    1 & 1\\
    0 & \\
  \end{tabular}}
  \end{minipage}\hfill
  \begin{minipage}[t]{0.24\linewidth}
  \textbf{d)}\\[2pt]
  {\small\renewcommand{\arraystretch}{1.0}\begin{tabular}{r|l}
    326 & 0\\
    163 & 1\\
    81 & 1\\
    40 & 0\\
    20 & 0\\
    10 & 0\\
    5 & 1\\
    2 & 0\\
    1 & 1\\
    0 & \\
  \end{tabular}}
  \end{minipage}\\[6pt]
\par\smallskip\noindent\hspace*{1em}\begin{tabular}{@{}p{0.42\linewidth}p{0.42\linewidth}@{}}
  \textbf{a)}~$489_{10} = 111101001_2$ & \textbf{b)}~$157_{10} = 10011101_2$\\[4pt]
  \textbf{c)}~$763_{10} = 1011111011_2$ & \textbf{d)}~$326_{10} = 101000110_2$
\end{tabular}\par\smallskip
  Si divide per 2 fino ad arrivare a 0 e si leggono i resti dal basso verso l'alto.}
  \criterioABC
  {2 pt per ogni numero: 1 pt divisioni per 2 corrette, 1 pt resti letti dal basso e risultato. Risultati: a) $10101101_2$; b) $111001101_2$; c) $101100110_2$; d) $1100111101_2$.}
  {2 pt per ogni numero: 1 pt divisioni per 2 corrette, 1 pt resti letti dal basso e risultato. Risultati: a) $110001010_2$; b) $11010011_2$; c) $1110010101_2$; d) $110110101_2$.}
  {2 pt per ogni numero: 1 pt divisioni per 2 corrette, 1 pt resti letti dal basso e risultato. Risultati: a) $111101001_2$; b) $10011101_2$; c) $1011111011_2$; d) $101000110_2$.}
\end{solution}
\end{esercizio}

% ---------------------------------------------------------------------
% Argomento: conversione binario -> decimale (somma delle potenze di 2)
%   4 numeri per fila (8, 8, 9 e 10 bit), un solo numero pari
%   fila A: 10110111=183, 101100101=357, 11010110=214, 1011000101=709
%   fila B: 110010011=403, 10100111=167, 1010100101=677, 11101100=236
%   fila C: 11001101=205, 1010001101=653, 101001011=331, 10110010=178
% Livello: base   |   Tempo stimato: 10 min
% ---------------------------------------------------------------------
\begin{esercizio}[8]
Converti in decimale:
\ABC{%
\par\smallskip\noindent\hspace*{1em}\begin{tabular}{@{}p{0.42\linewidth}p{0.42\linewidth}@{}}
  \textbf{a)}~$10110111_{2} = \ldots\ldots\ldots_{10}$ & \textbf{b)}~$101100101_{2} = \ldots\ldots\ldots_{10}$\\[4pt]
  \textbf{c)}~$11010110_{2} = \ldots\ldots\ldots_{10}$ & \textbf{d)}~$1011000101_{2} = \ldots\ldots\ldots_{10}$
\end{tabular}}{%
\par\smallskip\noindent\hspace*{1em}\begin{tabular}{@{}p{0.42\linewidth}p{0.42\linewidth}@{}}
  \textbf{a)}~$110010011_{2} = \ldots\ldots\ldots_{10}$ & \textbf{b)}~$10100111_{2} = \ldots\ldots\ldots_{10}$\\[4pt]
  \textbf{c)}~$1010100101_{2} = \ldots\ldots\ldots_{10}$ & \textbf{d)}~$11101100_{2} = \ldots\ldots\ldots_{10}$
\end{tabular}}{%
\par\smallskip\noindent\hspace*{1em}\begin{tabular}{@{}p{0.42\linewidth}p{0.42\linewidth}@{}}
  \textbf{a)}~$11001101_{2} = \ldots\ldots\ldots_{10}$ & \textbf{b)}~$1010001101_{2} = \ldots\ldots\ldots_{10}$\\[4pt]
  \textbf{c)}~$101001011_{2} = \ldots\ldots\ldots_{10}$ & \textbf{d)}~$10110010_{2} = \ldots\ldots\ldots_{10}$
\end{tabular}}

\begin{solution}
\ABC{%
  \textbf{a)}~$10110111_2 = 128 + 32 + 16 + 4 + 2 + 1 = 183_{10}$\\[4pt]
  \textbf{b)}~$101100101_2 = 256 + 64 + 32 + 4 + 1 = 357_{10}$\\[4pt]
  \textbf{c)}~$11010110_2 = 128 + 64 + 16 + 4 + 2 = 214_{10}$\\[4pt]
  \textbf{d)}~$1011000101_2 = 512 + 128 + 64 + 4 + 1 = 709_{10}$\\[4pt]
  Si sommano i pesi $2^0, 2^1, 2^2, \ldots$ (da destra) dei soli bit a 1.}{%
  \textbf{a)}~$110010011_2 = 256 + 128 + 16 + 2 + 1 = 403_{10}$\\[4pt]
  \textbf{b)}~$10100111_2 = 128 + 32 + 4 + 2 + 1 = 167_{10}$\\[4pt]
  \textbf{c)}~$1010100101_2 = 512 + 128 + 32 + 4 + 1 = 677_{10}$\\[4pt]
  \textbf{d)}~$11101100_2 = 128 + 64 + 32 + 8 + 4 = 236_{10}$\\[4pt]
  Si sommano i pesi $2^0, 2^1, 2^2, \ldots$ (da destra) dei soli bit a 1.}{%
  \textbf{a)}~$11001101_2 = 128 + 64 + 8 + 4 + 1 = 205_{10}$\\[4pt]
  \textbf{b)}~$1010001101_2 = 512 + 128 + 8 + 4 + 1 = 653_{10}$\\[4pt]
  \textbf{c)}~$101001011_2 = 256 + 64 + 8 + 2 + 1 = 331_{10}$\\[4pt]
  \textbf{d)}~$10110010_2 = 128 + 32 + 16 + 2 = 178_{10}$\\[4pt]
  Si sommano i pesi $2^0, 2^1, 2^2, \ldots$ (da destra) dei soli bit a 1.}
  \criterioABC
  {2 pt per ogni numero: 1 pt pesi (potenze di 2) corretti, 1 pt somma e risultato. Risultati: a) $183_{10}$; b) $357_{10}$; c) $214_{10}$; d) $709_{10}$.}
  {2 pt per ogni numero: 1 pt pesi (potenze di 2) corretti, 1 pt somma e risultato. Risultati: a) $403_{10}$; b) $167_{10}$; c) $677_{10}$; d) $236_{10}$.}
  {2 pt per ogni numero: 1 pt pesi (potenze di 2) corretti, 1 pt somma e risultato. Risultati: a) $205_{10}$; b) $653_{10}$; c) $331_{10}$; d) $178_{10}$.}
\end{solution}
\end{esercizio}

% ---------------------------------------------------------------------
% Argomento: 3 ingressi, 2 porte AND/OR in cascata, un ingresso negato (pallino) sulla prima porta
%   fila A: X = AND(A, NOT B),  Y = OR(X, C)
%   fila B: X = OR(NOT A, B),  Y = AND(X, C)
%   fila C: X = OR(B, NOT C),  Y = AND(X, A)
% Solo tabella di verita' (le espressioni logiche non sono ancora state spiegate)
% Livello: intermedio   |   Tempo stimato: 8 min
% ---------------------------------------------------------------------
\begin{esercizio}[6]
Scrivi la tabella di verità del circuito in figura (colonne $A$, $B$, $C$, $X$, $Y$).

\begin{center}
\ABC{%
  \begin{circuitikz}[x=1cm,y=1cm,scale=1.25,transform shape]
    \node[american and port] (g1) at (0,0.3) {};
    \node[american or port] (g2) at (3,-0.25) {};
    \draw (g1.in 1) -- ++(-0.9,0) node[left] {$A$};
    \draw (g1.in 2) -- ++(-0.9,0) node[left] {$B$};
    \draw (g2.in 2) -- (g1.in 2 |- g2.in 2) -- ++(-0.9,0) node[left] {$C$};
    \draw (g1.out) -- ++(0.45,0) node[above] {$X$} |- (g2.in 1);
    \draw (g2.out) -- ++(0.7,0) node[right] {$Y$};
    \node[circle,draw,fill=white,inner sep=0pt,minimum size=1.6mm,anchor=east] at (g1.bin 2) {};
  \end{circuitikz}}{%
  \begin{circuitikz}[x=1cm,y=1cm,scale=1.25,transform shape]
    \node[american or port] (g1) at (0,0.3) {};
    \node[american and port] (g2) at (3,-0.25) {};
    \draw (g1.in 1) -- ++(-0.9,0) node[left] {$A$};
    \draw (g1.in 2) -- ++(-0.9,0) node[left] {$B$};
    \draw (g2.in 2) -- (g1.in 2 |- g2.in 2) -- ++(-0.9,0) node[left] {$C$};
    \draw (g1.out) -- ++(0.45,0) node[above] {$X$} |- (g2.in 1);
    \draw (g2.out) -- ++(0.7,0) node[right] {$Y$};
    \node[circle,draw,fill=white,inner sep=0pt,minimum size=1.6mm,anchor=east] at (g1.bin 1) {};
  \end{circuitikz}}{%
  \begin{circuitikz}[x=1cm,y=1cm,scale=1.25,transform shape]
    \node[american or port] (g1) at (0,-0.3) {};
    \node[american and port] (g2) at (3,0.25) {};
    \draw (g1.in 1) -- ++(-0.9,0) node[left] {$B$};
    \draw (g1.in 2) -- ++(-0.9,0) node[left] {$C$};
    \draw (g2.in 1) -- (g1.in 1 |- g2.in 1) -- ++(-0.9,0) node[left] {$A$};
    \draw (g1.out) -- ++(0.45,0) node[below] {$X$} |- (g2.in 2);
    \draw (g2.out) -- ++(0.7,0) node[right] {$Y$};
    \node[circle,draw,fill=white,inner sep=0pt,minimum size=1.6mm,anchor=east] at (g1.bin 2) {};
  \end{circuitikz}}
\end{center}

\begin{solution}
\ABC{%
  {\renewcommand{\arraystretch}{1.05}\small
  \begin{tabular}{ccc|c|c|c}
    $A$ & $B$ & $C$ & NOT $B$ & $X$ & $Y$\\\hline
    0 & 0 & 0 & 1 & 0 & 0\\
    0 & 0 & 1 & 1 & 0 & 1\\
    0 & 1 & 0 & 0 & 0 & 0\\
    0 & 1 & 1 & 0 & 0 & 1\\
    1 & 0 & 0 & 1 & 1 & 1\\
    1 & 0 & 1 & 1 & 1 & 1\\
    1 & 1 & 0 & 0 & 0 & 0\\
    1 & 1 & 1 & 0 & 0 & 1\\
  \end{tabular}}}{%
  {\renewcommand{\arraystretch}{1.05}\small
  \begin{tabular}{ccc|c|c|c}
    $A$ & $B$ & $C$ & NOT $A$ & $X$ & $Y$\\\hline
    0 & 0 & 0 & 1 & 1 & 0\\
    0 & 0 & 1 & 1 & 1 & 1\\
    0 & 1 & 0 & 1 & 1 & 0\\
    0 & 1 & 1 & 1 & 1 & 1\\
    1 & 0 & 0 & 0 & 0 & 0\\
    1 & 0 & 1 & 0 & 0 & 0\\
    1 & 1 & 0 & 0 & 1 & 0\\
    1 & 1 & 1 & 0 & 1 & 1\\
  \end{tabular}}}{%
  {\renewcommand{\arraystretch}{1.05}\small
  \begin{tabular}{ccc|c|c|c}
    $A$ & $B$ & $C$ & NOT $C$ & $X$ & $Y$\\\hline
    0 & 0 & 0 & 1 & 1 & 0\\
    0 & 0 & 1 & 0 & 0 & 0\\
    0 & 1 & 0 & 1 & 1 & 0\\
    0 & 1 & 1 & 0 & 1 & 0\\
    1 & 0 & 0 & 1 & 1 & 1\\
    1 & 0 & 1 & 0 & 0 & 0\\
    1 & 1 & 0 & 1 & 1 & 1\\
    1 & 1 & 1 & 0 & 1 & 1\\
  \end{tabular}}}\\[4pt]
  Il pallino all'ingresso di una porta nega quel segnale prima della porta; il pallino all'uscita
  (NAND, NOR) nega il risultato della porta.
  \criterioABC
  {Le 8 combinazioni di $A$, $B$, $C$ tutte presenti: 1 pt; colonna $X$ (attenzione a NOT $B$): 2 pt; colonna $Y$: 3 pt ($-0{,}5$ per ogni riga sbagliata; un errore su $X$ non si ripenalizza su $Y$ se $Y$ è coerente). Colonne attese da 000 a 111: $X$ = 00001100, $Y$ = 01011101.}
  {Le 8 combinazioni di $A$, $B$, $C$ tutte presenti: 1 pt; colonna $X$ (attenzione a NOT $A$): 2 pt; colonna $Y$: 3 pt ($-0{,}5$ per ogni riga sbagliata; un errore su $X$ non si ripenalizza su $Y$ se $Y$ è coerente). Colonne attese da 000 a 111: $X$ = 11110011, $Y$ = 01010001.}
  {Le 8 combinazioni di $A$, $B$, $C$ tutte presenti: 1 pt; colonna $X$ (attenzione a NOT $C$): 2 pt; colonna $Y$: 3 pt ($-0{,}5$ per ogni riga sbagliata; un errore su $X$ non si ripenalizza su $Y$ se $Y$ è coerente). Colonne attese da 000 a 111: $X$ = 10111011, $Y$ = 00001011.}
\end{solution}
\end{esercizio}

% ---------------------------------------------------------------------
% Argomento: 3 ingressi, 2 porte in cascata con NAND/NOR e un ingresso negato sulla seconda porta
%   fila A: X = OR(A, B),  Y = NAND(X, NOT C)
%   fila B: X = AND(A, B),  Y = NOR(X, NOT C)
%   fila C: X = NOR(B, C),  Y = AND(X, NOT A)
% Solo tabella di verita' (le espressioni logiche non sono ancora state spiegate)
% Livello: intermedio   |   Tempo stimato: 8 min
% ---------------------------------------------------------------------
\begin{esercizio}[6]
Scrivi la tabella di verità del circuito in figura.

\begin{center}
\ABC{%
  \begin{circuitikz}[x=1cm,y=1cm,scale=1.25,transform shape]
    \node[american or port] (g1) at (0,0.3) {};
    \node[american nand port] (g2) at (3,-0.25) {};
    \draw (g1.in 1) -- ++(-0.9,0) node[left] {$A$};
    \draw (g1.in 2) -- ++(-0.9,0) node[left] {$B$};
    \draw (g2.in 2) -- (g1.in 2 |- g2.in 2) -- ++(-0.9,0) node[left] {$C$};
    \draw (g1.out) -- ++(0.45,0) |- (g2.in 1);
    \draw (g2.out) -- ++(0.7,0) node[right] {$Y$};
    \node[circle,draw,fill=white,inner sep=0pt,minimum size=1.6mm,anchor=east] at (g2.bin 2) {};
  \end{circuitikz}}{%
  \begin{circuitikz}[x=1cm,y=1cm,scale=1.25,transform shape]
    \node[american and port] (g1) at (0,0.3) {};
    \node[american nor port] (g2) at (3,-0.25) {};
    \draw (g1.in 1) -- ++(-0.9,0) node[left] {$A$};
    \draw (g1.in 2) -- ++(-0.9,0) node[left] {$B$};
    \draw (g2.in 2) -- (g1.in 2 |- g2.in 2) -- ++(-0.9,0) node[left] {$C$};
    \draw (g1.out) -- ++(0.45,0) |- (g2.in 1);
    \draw (g2.out) -- ++(0.7,0) node[right] {$Y$};
    \node[circle,draw,fill=white,inner sep=0pt,minimum size=1.6mm,anchor=east] at (g2.bin 2) {};
  \end{circuitikz}}{%
  \begin{circuitikz}[x=1cm,y=1cm,scale=1.25,transform shape]
    \node[american nor port] (g1) at (0,-0.3) {};
    \node[american and port] (g2) at (3,0.25) {};
    \draw (g1.in 1) -- ++(-0.9,0) node[left] {$B$};
    \draw (g1.in 2) -- ++(-0.9,0) node[left] {$C$};
    \draw (g2.in 1) -- (g1.in 1 |- g2.in 1) -- ++(-0.9,0) node[left] {$A$};
    \draw (g1.out) -- ++(0.45,0) |- (g2.in 2);
    \draw (g2.out) -- ++(0.7,0) node[right] {$Y$};
    \node[circle,draw,fill=white,inner sep=0pt,minimum size=1.6mm,anchor=east] at (g2.bin 1) {};
  \end{circuitikz}}
\end{center}

\begin{solution}
\ABC{%
\begin{center}
  \begin{circuitikz}[x=1cm,y=1cm,scale=1.25,transform shape]
    \node[american or port] (g1) at (0,0.3) {};
    \node[american nand port] (g2) at (3,-0.25) {};
    \draw (g1.in 1) -- ++(-0.9,0) node[left] {$A$};
    \draw (g1.in 2) -- ++(-0.9,0) node[left] {$B$};
    \draw (g2.in 2) -- (g1.in 2 |- g2.in 2) -- ++(-0.9,0) node[left] {$C$};
    \draw (g1.out) -- ++(0.45,0) |- (g2.in 1);
    \draw (g2.out) -- ++(0.7,0) node[right] {$Y$};
    \node[circle,draw,fill=white,inner sep=0pt,minimum size=1.6mm,anchor=east] at (g2.bin 2) {};
  \end{circuitikz}
\end{center}
  {\renewcommand{\arraystretch}{1.05}\small
  \begin{tabular}{ccc|c}
    $A$ & $B$ & $C$ & $Y$\\\hline
    0 & 0 & 0 & 1\\
    0 & 0 & 1 & 1\\
    0 & 1 & 0 & 0\\
    0 & 1 & 1 & 1\\
    1 & 0 & 0 & 0\\
    1 & 0 & 1 & 1\\
    1 & 1 & 0 & 0\\
    1 & 1 & 1 & 1\\
  \end{tabular}}}{%
\begin{center}
  \begin{circuitikz}[x=1cm,y=1cm,scale=1.25,transform shape]
    \node[american and port] (g1) at (0,0.3) {};
    \node[american nor port] (g2) at (3,-0.25) {};
    \draw (g1.in 1) -- ++(-0.9,0) node[left] {$A$};
    \draw (g1.in 2) -- ++(-0.9,0) node[left] {$B$};
    \draw (g2.in 2) -- (g1.in 2 |- g2.in 2) -- ++(-0.9,0) node[left] {$C$};
    \draw (g1.out) -- ++(0.45,0) |- (g2.in 1);
    \draw (g2.out) -- ++(0.7,0) node[right] {$Y$};
    \node[circle,draw,fill=white,inner sep=0pt,minimum size=1.6mm,anchor=east] at (g2.bin 2) {};
  \end{circuitikz}
\end{center}
  {\renewcommand{\arraystretch}{1.05}\small
  \begin{tabular}{ccc|c}
    $A$ & $B$ & $C$ & $Y$\\\hline
    0 & 0 & 0 & 0\\
    0 & 0 & 1 & 1\\
    0 & 1 & 0 & 0\\
    0 & 1 & 1 & 1\\
    1 & 0 & 0 & 0\\
    1 & 0 & 1 & 1\\
    1 & 1 & 0 & 0\\
    1 & 1 & 1 & 0\\
  \end{tabular}}}{%
\begin{center}
  \begin{circuitikz}[x=1cm,y=1cm,scale=1.25,transform shape]
    \node[american nor port] (g1) at (0,-0.3) {};
    \node[american and port] (g2) at (3,0.25) {};
    \draw (g1.in 1) -- ++(-0.9,0) node[left] {$B$};
    \draw (g1.in 2) -- ++(-0.9,0) node[left] {$C$};
    \draw (g2.in 1) -- (g1.in 1 |- g2.in 1) -- ++(-0.9,0) node[left] {$A$};
    \draw (g1.out) -- ++(0.45,0) |- (g2.in 2);
    \draw (g2.out) -- ++(0.7,0) node[right] {$Y$};
    \node[circle,draw,fill=white,inner sep=0pt,minimum size=1.6mm,anchor=east] at (g2.bin 1) {};
  \end{circuitikz}
\end{center}
  {\renewcommand{\arraystretch}{1.05}\small
  \begin{tabular}{ccc|c}
    $A$ & $B$ & $C$ & $Y$\\\hline
    0 & 0 & 0 & 1\\
    0 & 0 & 1 & 0\\
    0 & 1 & 0 & 0\\
    0 & 1 & 1 & 0\\
    1 & 0 & 0 & 0\\
    1 & 0 & 1 & 0\\
    1 & 1 & 0 & 0\\
    1 & 1 & 1 & 0\\
  \end{tabular}}}
  \criterioABC
  {Le 8 combinazioni di $A$, $B$, $C$ tutte presenti: 1 pt; colonna $Y$ corretta: 5 pt ($-1$ per ogni riga sbagliata, minimo 0; attenzione a NOT $C$). Colonna $Y$ attesa da 000 a 111: 11010101 (uscita della prima porta, per controllare i passaggi: 00111111).}
  {Le 8 combinazioni di $A$, $B$, $C$ tutte presenti: 1 pt; colonna $Y$ corretta: 5 pt ($-1$ per ogni riga sbagliata, minimo 0; attenzione a NOT $C$). Colonna $Y$ attesa da 000 a 111: 01010100 (uscita della prima porta, per controllare i passaggi: 00000011).}
  {Le 8 combinazioni di $A$, $B$, $C$ tutte presenti: 1 pt; colonna $Y$ corretta: 5 pt ($-1$ per ogni riga sbagliata, minimo 0; attenzione a NOT $A$). Colonna $Y$ attesa da 000 a 111: 10000000 (uscita della prima porta, per controllare i passaggi: 10001000).}
\end{solution}
\end{esercizio}

\sceltaesercizi{Diagrammi di flusso: svolgi a scelta UNO SOLO tra l'esercizio 5 e l'esercizio 6.
  Con l'esercizio 5 (più semplice) il punteggio massimo della verifica è 34 punti, con
  l'esercizio 6 è 42 punti.}
% ---------------------------------------------------------------------
% Argomento: diagramma di flusso con calcolo e due decisioni annidate
%   fila A: forno, media di due temperature -> resistenza / ventola / ok
%   fila B: serbatoio, percentuale di riempimento -> carico / scarico / ok
%   fila C: totale di un acquisto con sconto del 20 % o del 10 %
% Livello: intermedio   |   Tempo stimato: 15 min
% ---------------------------------------------------------------------
\begin{esercizio}[8]
Disegna il diagramma di flusso di un algoritmo che
\ABC{%
legge le temperature $T_1$ e $T_2$ misurate da due sensori in un forno e ne calcola la
media $T$. Se $T$ è minore di \qty{180}{\degreeCelsius} visualizza ``Accendi la resistenza'', se è
maggiore di \qty{220}{\degreeCelsius} visualizza ``Accendi la ventola'', altrimenti visualizza
``Temperatura corretta''.}{%
legge il livello $L$ dell'acqua in un serbatoio e l'altezza $H$ del serbatoio (entrambi in cm)
e calcola la percentuale di riempimento $P$. Se $P$ è minore del \qty{20}{\percent} visualizza
``Apri il carico'', se è maggiore del \qty{90}{\percent} visualizza ``Apri lo scarico'',
altrimenti visualizza ``Livello corretto''.}{%
legge il prezzo $P$ di un prodotto e la quantità $Q$ acquistata e calcola il totale $T$.
Se il totale supera \qty{100}{\eur} applica uno sconto del \qty{20}{\percent}, se supera
\qty{50}{\eur} (ma non \qty{100}{\eur}) uno sconto del \qty{10}{\percent}, altrimenti nessuno
sconto. Alla fine visualizza il totale da pagare.}

\begin{solution}
\begin{center}
\ABC{%
\begin{tikzpicture}[x=1cm,y=1cm]
  \node[fcterm] (ini) at (0,0) {Inizio};
  \node[fcio]   (in)  at (0,-1.1) {Leggi $T_1$, $T_2$};
  \node[fcproc] (med) at (0,-2.2) {$T \gets (T_1+T_2)/2$};
  \node[fcdec]  (d1)  at (0,-3.55) {$T < 180$};
  \node[fcio]   (o1)  at (4.6,-3.55) {Scrivi ``Accendi la resistenza''};
  \node[fcdec]  (d2)  at (0,-5.25) {$T > 220$};
  \node[fcio]   (o2)  at (4.6,-5.25) {Scrivi ``Accendi la ventola''};
  \node[fcio]   (o3)  at (0,-6.75) {Scrivi ``Temperatura corretta''};
  \node[fcterm] (fin) at (0,-8.4) {Fine};
  \draw[fcfreccia] (ini) -- (in);
  \draw[fcfreccia] (in) -- (med);
  \draw[fcfreccia] (med) -- (d1);
  \draw[fcfreccia] (d1) -- node[fcsi,above] {Sì} (o1);
  \draw[fcfreccia] (d1) -- node[fcno,right] {No} (d2);
  \draw[fcfreccia] (d2) -- node[fcsi,above] {Sì} (o2);
  \draw[fcfreccia] (d2) -- node[fcno,right] {No} (o3);
  \draw[thick] (o1.east) -- (7.4,-3.55) |- (0,-7.65);
  \draw[thick] (o2.east) -- (7.4,-5.25);
  \draw[fcfreccia] (o3) -- (fin);
\end{tikzpicture}}{%
\begin{tikzpicture}[x=1cm,y=1cm]
  \node[fcterm] (ini) at (0,0) {Inizio};
  \node[fcio]   (in)  at (0,-1.1) {Leggi $L$, $H$};
  \node[fcproc] (med) at (0,-2.2) {$P \gets L/H \cdot 100$};
  \node[fcdec]  (d1)  at (0,-3.55) {$P < 20$};
  \node[fcio]   (o1)  at (4.4,-3.55) {Scrivi ``Apri il carico''};
  \node[fcdec]  (d2)  at (0,-5.25) {$P > 90$};
  \node[fcio]   (o2)  at (4.4,-5.25) {Scrivi ``Apri lo scarico''};
  \node[fcio]   (o3)  at (0,-6.75) {Scrivi ``Livello corretto''};
  \node[fcterm] (fin) at (0,-8.4) {Fine};
  \draw[fcfreccia] (ini) -- (in);
  \draw[fcfreccia] (in) -- (med);
  \draw[fcfreccia] (med) -- (d1);
  \draw[fcfreccia] (d1) -- node[fcsi,above] {Sì} (o1);
  \draw[fcfreccia] (d1) -- node[fcno,right] {No} (d2);
  \draw[fcfreccia] (d2) -- node[fcsi,above] {Sì} (o2);
  \draw[fcfreccia] (d2) -- node[fcno,right] {No} (o3);
  \draw[thick] (o1.east) -- (6.8,-3.55) |- (0,-7.65);
  \draw[thick] (o2.east) -- (6.8,-5.25);
  \draw[fcfreccia] (o3) -- (fin);
\end{tikzpicture}}{%
\begin{tikzpicture}[x=1cm,y=1cm]
  \node[fcterm] (ini) at (0,0) {Inizio};
  \node[fcio]   (in)  at (0,-1.1) {Leggi $P$, $Q$};
  \node[fcproc] (tot) at (0,-2.2) {$T \gets P \cdot Q$};
  \node[fcdec]  (d1)  at (0,-3.55) {$T > 100$};
  \node[fcproc] (s1)  at (3.8,-3.55) {$T \gets T \cdot 0{,}8$};
  \node[fcdec]  (d2)  at (0,-5.25) {$T > 50$};
  \node[fcproc] (s2)  at (3.8,-5.25) {$T \gets T \cdot 0{,}9$};
  \node[fcio]   (out) at (0,-7.15) {Scrivi $T$};
  \node[fcterm] (fin) at (0,-8.4) {Fine};
  \draw[fcfreccia] (ini) -- (in);
  \draw[fcfreccia] (in) -- (tot);
  \draw[fcfreccia] (tot) -- (d1);
  \draw[fcfreccia] (d1) -- node[fcsi,above] {Sì} (s1);
  \draw[fcfreccia] (d1) -- node[fcno,right] {No} (d2);
  \draw[fcfreccia] (d2) -- node[fcsi,above] {Sì} (s2);
  \draw[fcfreccia] (d2) -- node[fcno,right] {No} (out);
  \draw[thick] (s1.east) -- (5.4,-3.55) |- (0,-6.35);
  \draw[thick] (s2.east) -- (5.4,-5.25);
  \draw[fcfreccia] (out) -- (fin);
\end{tikzpicture}}
\end{center}
\ABC{%
  Il secondo controllo si fa solo se $T \geq 180$: quindi nel ramo ``No'' vale già $T \geq 180$ e
  basta controllare $T > 220$.}{%
  Il secondo controllo si fa solo se $P \geq 20$: quindi nel ramo ``No'' vale già $P \geq 20$ e
  basta controllare $P > 90$.}{%
  Il secondo controllo si fa solo se $T \leq 100$: quindi nel ramo ``No'' basta controllare $T > 50$.
  Sconto del \qty{20}{\percent}: si paga l'\qty{80}{\percent}, cioè $T \cdot 0{,}8$
  (oppure $T \gets T - T\cdot 20/100$).}
  \criterioABC
  {Inizio/Fine e simboli corretti (ovale, parallelogramma, rettangolo, rombo): 1 pt; lettura di $T_1$ e
    $T_2$: 1 pt; calcolo della media $(T_1+T_2)/2$: 1 pt (senza parentesi: 0 pt); prima decisione
    $T<180$ con i rami Sì/No: 2 pt; seconda decisione $T>220$ nel ramo No: 2 pt; tre messaggi
    nei rami giusti e rami che si ricongiungono prima di Fine: 1 pt.}
  {Inizio/Fine e simboli corretti (ovale, parallelogramma, rettangolo, rombo): 1 pt; lettura di $L$ e
    $H$: 1 pt; calcolo $P = L/H\cdot 100$: 1 pt; prima decisione $P<20$ con i rami Sì/No: 2 pt;
    seconda decisione $P>90$ nel ramo No: 2 pt; tre messaggi nei rami giusti e rami che si
    ricongiungono prima di Fine: 1 pt.}
  {Inizio/Fine e simboli corretti (ovale, parallelogramma, rettangolo, rombo): 1 pt; lettura di $P$ e
    $Q$: 1 pt; calcolo $T = P\cdot Q$: 1 pt; prima decisione $T>100$ con sconto del 20\,\%: 2 pt;
    seconda decisione $T>50$ nel ramo No con sconto del 10\,\%: 2 pt; rami che si ricongiungono e
    una sola uscita ``Scrivi $T$'' (o una per ramo): 1 pt.}
\end{solution}
\end{esercizio}

% ---------------------------------------------------------------------
% Argomento: diagramma di flusso con ciclo (niente vettori: un solo dato
%            letto, il ciclo aggiorna variabili con nome dato nel testo)
%   fila A: serbatoio riempito a secchi -> numero di secchi e livello finale
%   fila B: magazzino svuotato a spedizioni -> numero di spedizioni e pezzi rimasti
%   (niente tempo: il ciclo conta operazioni, non minuti)
%   fila C: scaffale riempito con una pila di scatole -> numero di scatole e altezza della pila
% Livello: intermedio (nessun suggerimento: condizione e aggiornamenti da ricavare)   |   Tempo stimato: 15 min
% ---------------------------------------------------------------------
\begin{esercizio}[12]
\ABC{%
Un serbatoio, inizialmente vuoto, si riempie versando un secchio d'acqua alla volta.
Disegna il diagramma di flusso di un algoritmo che conta quanti secchi servono per riempirlo.\\
\textbf{Dati da leggere:} $C$ = capacità del serbatoio (in litri); $S$ = litri contenuti in un secchio.\\
\textbf{Altre variabili:} $L$ = litri già presenti nel serbatoio; $N$ = numero di secchi versati.\\
Si versa un secchio dopo l'altro finché $L$ non raggiunge la capacità $C$. Alla fine visualizza $N$ e $L$.}{%
Un magazzino si svuota con una spedizione dopo l'altra. Disegna il diagramma di flusso di un
algoritmo che conta quante spedizioni si possono fare.\\
\textbf{Dati da leggere:} $Q$ = pezzi presenti in magazzino; $S$ = pezzi che partono con ogni spedizione.\\
\textbf{Altre variabili:} $K$ = numero di spedizioni fatte.\\
Si fanno spedizioni finché in magazzino ci sono abbastanza pezzi per una spedizione completa.
Alla fine visualizza $K$ e i pezzi rimasti $Q$.}{%
Su uno scaffale si impilano scatole tutte uguali, una sopra l'altra. Disegna il diagramma di
flusso di un algoritmo che conta quante scatole ci stanno.\\
\textbf{Dati da leggere:} $A$ = altezza dello scaffale (in cm); $H$ = altezza di una scatola (in cm).\\
\textbf{Altre variabili:} $P$ = altezza della pila; $N$ = numero di scatole impilate.\\
Si aggiungono scatole finché ce ne sta ancora una: la pila non deve mai superare l'altezza $A$.
Alla fine visualizza $N$ e $P$.}

\begin{solution}
\begin{center}
\ABC{%
\begin{tikzpicture}[x=1cm,y=1cm]
  \node[fcterm] (ini)  at (0,0) {Inizio};
  \node[fcio]   (in)   at (0,-1.1) {Leggi $C$, $S$};
  \node[fcinit] (ini0) at (0,-2.2) {$L \gets 0$, \ $N \gets 0$};
  \node[fcdec]  (d1)   at (0,-3.75) {$L < C$};
  \node[fcproc] (agg)  at (0,-5.3) {$L \gets L + S$};
  \node[fcproc] (inc)  at (0,-6.4) {$N \gets N + 1$};
  \node[fcio]   (out)  at (5.0,-3.75) {Scrivi $N$, $L$};
  \node[fcterm] (fin)  at (5.0,-4.95) {Fine};
  \draw[fcfreccia] (ini) -- (in);
  \draw[fcfreccia] (in) -- (ini0);
  \draw[fcfreccia] (ini0) -- (d1);
  \draw[fcfreccia] (d1) -- node[fcsi,right] {Sì} (agg);
  \draw[fcfreccia] (agg) -- (inc);
  \draw[fcfreccia] (inc.south) -- ++(0,-0.5) -| (-2.4,-2.95) -- (0,-2.95);
  \draw[fcfreccia] (d1) -- node[fcno,above] {No} (out);
  \draw[fcfreccia] (out) -- (fin);
\end{tikzpicture}}{%
\begin{tikzpicture}[x=1cm,y=1cm]
  \node[fcterm] (ini)  at (0,0) {Inizio};
  \node[fcio]   (in)   at (0,-1.1) {Leggi $Q$, $S$};
  \node[fcinit] (ini0) at (0,-2.2) {$K \gets 0$};
  \node[fcdec]  (d1)   at (0,-3.75) {$Q \geq S$};
  \node[fcproc] (agg)  at (0,-5.3) {$Q \gets Q - S$};
  \node[fcproc] (inc)  at (0,-6.4) {$K \gets K + 1$};
  \node[fcio]   (out)  at (5.0,-3.75) {Scrivi $K$, $Q$};
  \node[fcterm] (fin)  at (5.0,-4.95) {Fine};
  \draw[fcfreccia] (ini) -- (in);
  \draw[fcfreccia] (in) -- (ini0);
  \draw[fcfreccia] (ini0) -- (d1);
  \draw[fcfreccia] (d1) -- node[fcsi,right] {Sì} (agg);
  \draw[fcfreccia] (agg) -- (inc);
  \draw[fcfreccia] (inc.south) -- ++(0,-0.5) -| (-2.4,-2.95) -- (0,-2.95);
  \draw[fcfreccia] (d1) -- node[fcno,above] {No} (out);
  \draw[fcfreccia] (out) -- (fin);
\end{tikzpicture}}{%
\begin{tikzpicture}[x=1cm,y=1cm]
  \node[fcterm] (ini)  at (0,0) {Inizio};
  \node[fcio]   (in)   at (0,-1.1) {Leggi $A$, $H$};
  \node[fcinit] (ini0) at (0,-2.2) {$P \gets 0$, \ $N \gets 0$};
  \node[fcdec]  (d1)   at (0,-3.75) {$P + H \leq A$};
  \node[fcproc] (agg)  at (0,-5.3) {$P \gets P + H$};
  \node[fcproc] (inc)  at (0,-6.4) {$N \gets N + 1$};
  \node[fcio]   (out)  at (5.0,-3.75) {Scrivi $N$, $P$};
  \node[fcterm] (fin)  at (5.0,-4.95) {Fine};
  \draw[fcfreccia] (ini) -- (in);
  \draw[fcfreccia] (in) -- (ini0);
  \draw[fcfreccia] (ini0) -- (d1);
  \draw[fcfreccia] (d1) -- node[fcsi,right] {Sì} (agg);
  \draw[fcfreccia] (agg) -- (inc);
  \draw[fcfreccia] (inc.south) -- ++(0,-0.5) -| (-2.4,-2.95) -- (0,-2.95);
  \draw[fcfreccia] (d1) -- node[fcno,above] {No} (out);
  \draw[fcfreccia] (out) -- (fin);
\end{tikzpicture}}
\end{center}
\ABC{%
  Ogni giro del ciclo è un secchio versato: $N$ conta i giri. $L$ parte da 0 (serbatoio vuoto);
  quando il ciclo finisce, $L$ può superare di poco $C$ perché aumenta di $S$ alla volta.}{%
  Ogni giro del ciclo è una spedizione: $K$ conta i giri. Quando il ciclo finisce, in magazzino
  restano meno di $S$ pezzi (possono essere anche 0).}{%
  Ogni giro del ciclo è una scatola aggiunta: $N$ conta i giri. Prima di aggiungere una scatola si
  controlla che la pila con la scatola in più ($P+H$) non superi $A$; alla fine $P \leq A$.}
  \criterioABC
  {Inizio/Fine e simboli corretti (inizializzazione nell'esagono): 1 pt; lettura di $C$ e $S$: 1 pt;
    inizializzazione $L\gets 0$, $N\gets 0$: 2 pt; condizione del ciclo $L<C$ ricavata dal testo, con
    rami Sì/No giusti: 3 pt; aggiornamento $L\gets L+S$: 2 pt; contatore $N\gets N+1$: 1 pt; ritorno alla
    condizione: 1 pt; uscita di $N$ e $L$ dopo il ciclo: 1 pt. Ciclo con condizione in fondo
    (ripeti\ldots finché $L \geq C$) corretto: stessi punti.}
  {Inizio/Fine e simboli corretti (inizializzazione nell'esagono): 1 pt; lettura di $Q$ e $S$: 1 pt;
    inizializzazione $K\gets 0$ (e nessuna inizializzazione di $Q$ dopo la lettura): 2 pt; condizione del
    ciclo $Q\geq S$ ricavata da ``abbastanza pezzi per una spedizione completa'', con rami Sì/No giusti:
    3 pt ($Q>S$, che perde la spedizione quando $Q=S$: 1 pt); aggiornamento $Q\gets Q-S$: 2 pt; contatore $K\gets K+1$: 1 pt; ritorno
    alla condizione: 1 pt; uscita di $K$ e $Q$ dopo il ciclo: 1 pt. Ciclo con condizione in fondo
    corretto se gestisce il caso $Q<S$ iniziale: stessi punti.}
  {Inizio/Fine e simboli corretti (inizializzazione nell'esagono): 1 pt; lettura di $A$ e $H$: 1 pt;
    inizializzazione $P\gets 0$, $N\gets 0$: 2 pt; condizione del ciclo $P+H\leq A$ ricavata da ``ce ne
    sta ancora una'', con rami Sì/No giusti: 3 pt ($P<A$ o $P\leq A$, che fanno superare $A$: 1 pt);
    aggiornamento $P\gets P+H$: 2 pt; contatore $N\gets N+1$: 1 pt; ritorno alla condizione: 1 pt; uscita
    di $N$ e $P$ dopo il ciclo: 1 pt. Ciclo con condizione in fondo corretto se gestisce il caso $H>A$:
    stessi punti.}
\end{solution}
\end{esercizio}


\end{questions}

% Tabella punti + griglia voti in fondo alla 2ª pagina
\riepilogo

\end{document}
